\section{Caption Hallucination Assessment}

We first introduce our image relevance metric, \emph{CHAIR}, which assesses captions w.r.t. objects that are actually in an image. It is used as a main tool in our evaluation. Next we discuss the notions of \emph{image and language model consistency}, which we use to reason about the causes of hallucination. %

\subsection{The \metric{} Metric}
\label{sec:metric}
%
To measure object hallucination, we propose the \emph{\metric{} (\metricfullname{})} metric, which calculates what proportion of words generated are actually in the image according to the ground truth sentences and object segmentations. This metric has two variants: per-instance, or what fraction of object instances are hallucinated (denoted as \metric{}i), and per-sentence, or what fraction of sentences include a hallucinated object (denoted as \metric{}s):%
$$\text{CHAIR}_{i} = \frac{ \vert\{ \text{hallucinated objects}\}\vert }{\vert\{\text{all objects mentioned}\}\vert }$$
$$\text{CHAIR}_{s} = \frac{ \vert\{ \text{sentences with hallucinated object}\}\vert }{\vert\{\text{  all sentences}\}\vert }$$

%
For easier analysis, we restrict our study to the 80 MSCOCO objects which appear in the MSCOCO segmentation challenge.
To determine whether a generated sentence contains hallucinated objects, we first tokenize each sentence and then singularize each word. 
We then use a list of synonyms for MSCOCO objects (based on the list from~\citet{lu2018neural}) to map words (e.g., ``player'') to MSCOCO objects (e.g., ``person''). 
Additionally, for sentences which include two word compounds (e.g., ``hot dog'') we take care that other MSCOCO objects (in this case ``dog'') are not incorrectly assigned to the list of MSCOCO objects in the sentence.
%
%
For each ground truth sentence, we determine a list of MSCOCO objects in the same way.
The MSCOCO segmentation annotations are used by simply relying on the provided object labels.

We find that considering both sources of annotation is important.
For example, MSCOCO contains an object ``dining table'' annotated with segmentation maps. However, humans refer to many different kinds of objects as ``table'' (e.g., ``coffee table'' or ``side table''), though these objects are not annotated as they are not specifically ``dining table''.
By using sentence annotations to scrape ground truth objects, we account for variation in how human annotators refer to different objects.
%
Inversely, we find that frequently humans will not mention all objects in a scene. 
Qualitatively, we observe that both annotations are important to capture hallucination.
Empirically, we verify that using only segmentation labels or only reference captions leads to higher hallucination (and practically incorrect) rates.
%

%
%

%
%

%
%
%

%


%

\subsection{Image Consistency}
\label{sec:img_consistency}
%
We define a notion of \textit{image consistency}, or how consistent errors from the captioning model are with a model which predicts objects based on an image alone.
To measure image consistency for a particular generated word, we train an image model and record $P(w|I)$ or the probability of predicting the word given only the image.
To score the image consistency of a caption we use the average of $P(w|I)$ for all MSCOCO objects, where
higher values mean that errors are \textit{more} consistent with the image model.
Our image model is a multi-label classification model with labels corresponding to MSCOCO objects (labels determined the same way as is done for \metric{}) which shares the visual features with the caption models.

\subsection{Language Consistency}
\label{sec:lang_consistency}
We also introduce a notion of \textit{language consistency}, i.e. how consistent errors from the captioning model are with a model which predicts words based only on previously generated words. We train an LSTM~\cite{hochreiter1997long} based language model which predicts a word $w_t$ given previous words $w_{0:t-1}$ 
%
 on MSCOCO data. 
We report language consistency as $1/R(w_t)$ where $R(w_t)$ is the rank of the predicted word in the language model.
Again, for a caption we report average rank across all MSCOCO objects in the sentence and higher language consistency implies that errors are \textit{more} consistent with the language model.

We illustrate image and language consistency in Figure~\ref{fig:lm_im}, i.e. the hallucination error (``fork'') is more consistent with the Language Model predictions than with the Image Model predictions. We use these consistency measures in Section \ref{sec:causes} to help us investigate the causes of hallucination.

\begin{figure}[t]
\centering
\includegraphics[width=\linewidth]{figures/LM_IM_example.pdf}
\caption{\label{fig:lm_im} Example of image and language consistency. The hallucination error (``fork'') is more consistent with the Language Model.}
\end{figure}

%
%
%
%

%

